% SPDX-License-Identifier: CC-BY-4.0
% =====================================================================
% Section 6 of D. Joyner, "A hitch-hiker's guide to the metaplectic trace
% formula for SL(2)", revised for the 2-fold cover (n = 2), November 2026.
%
% Stage 5 of ../revision_log.md. Sources: ../original-1995/cctrace6.tex
% (= ../book-1997/book7g.tex; ../book-1997/book7gnw.tex is a later version
% of the same chapter). Referee items II.14-II.17 and the Section 6 minor
% items are fixed here; see the "Stage 5 notes" in ../revision_log.md,
% which also map the 1995 numbers to the numbers here.
%
% This file is \input by the master; it compiles on its own through
% preview.tex. Explicit formulas are checked by checks/sec6_n2_check.py.
% Comments beginning "% [Claude]" are editorial notes by Claude.
% Revised after the second referee report (ChatGPT, November 2026): an
% explanation of the correction term before the formal construction, real
% normalizing factors for all unitary characters, the three senses of
% "normalizing factor", a fuller proof of M(w,0) = -1, and an accurate
% statement of what is explicit; comparisons with the 1995 notes moved to
% Appendix A.
% =====================================================================

\providecommand{\GA}{\overline{G(\A)}}
\providecommand{\GS}{\overline{G(\Q_S)}}
\providecommand{\AS}{\overline{A(\Q_S)}}
\providecommand{\KS}{\overline{K_S}}

\section{The invariant form}\label{sec:invariant}

The trace formula of Theorem~\ref{thm:tf} is an identity $\sum_{\mathcal O}
J_{\mathcal O}(f)=\sum_\chi J_\chi(f)$ between distributions, but not all
of the terms are invariant: $J(f^y)\ne J(f)$ in general, where
$f^y(x)=f(yxy^{-1})$. The elliptic terms $J_{\mathcal O}$ and the cuspidal
terms $J_\chi$ are invariant; the hyperbolic and central terms, and the
terms $J_\chi$ with $\chi\in\mathcal X_A$, are not. This section computes
the failure of invariance, defines the invariant distributions
$I_A(\gamma,f)$ that replace the weighted orbital integrals
$J_A(\gamma,f)$, following Arthur~\cite{A6,A9,A10}, and puts the trace
formula in invariant form (Theorem~\ref{thm:itf} and
Corollary~\ref{cor:itf-final}).

For covering groups the general theory is due to Li. The normalization
of the local intertwining operators, which the construction needs, is
in~\cite{LiII}. The invariant trace formula is in~\cite{LiIV}, under the
hypothesis that the invariant Paley--Wiener theorem holds at the real
places for the Levi subgroups; Li verifies the hypothesis for the
two-fold covers of $Sp(2n)$, hence for the group here. For
$G=\Res_{k/\Q}SL(2)$ much less is needed. The only proper Levi subgroup is
$A$, and its cover is abelian (Lemma~\ref{lem:global}(e)), so the one
Paley--Wiener theorem that enters is Fourier inversion on the abelian
group $\AS$. Given the normalized intertwining operators, we can
therefore prove the variance formulas, the invariance of $I_A(\gamma,f)$
and the invariant trace formula directly; we do so, and cite
Arthur and Li for the general statements.

\paragraph{Why a correction term, and why this one.} The weighted
orbital integral $J_A(\gamma,f)$ integrates $f(x^{-1}\gamma x)$ against the
weight $v_A(x)=-H(x)-H(wx)$. The height $H$ is defined through the
Iwasawa decomposition, so it depends on the choice of the maximal compact
subgroup: $H(xy)\ne H(x)+H(y)$ in general, but $H(xy)=H(x)+H(\tilde ky)$
when $x=n\tilde a\tilde k$. Replacing $f$ by $f^y$ replaces $v_A(x)$ by
$v_A(xy)$, and the difference $-H(\tilde ky)-H(\tilde k'y)$ does not
vanish. This is the geometric source of the non-invariance
(Lemma~\ref{lem:variance-local}). Since the difference depends only on the
$K$-components, the change in $J_A(\gamma,f)$ is not a new kind of
object: it is the value at $\gamma$ and at $w\gamma w^{-1}$ of one function
$f_{B,y}$ on the cover of $A$.

On the spectral side the same function appears. The logarithmic
derivative $\mathcal M^{-1}\mathcal M'$ of an intertwining operator
between $I(\pi_\lambda)$ and $I(\pi^w_{-\lambda})$ is not an intertwining
operator, because the two representations depend on $\lambda$ through
the same height: differentiating $e^{(\lambda+1)H(\tilde ky)}$ brings down
$H(\tilde ky)$. Conjugating $\mathcal M^{-1}\mathcal M'$ by $I(\pi_\lambda,
y)$ therefore changes it by multiplication by $H(\tilde ky)$, once
directly and once through $\mathcal M$ (Lemma~\ref{lem:conj-identity}),
and the trace against $I(\pi_\lambda,f)$ changes by the Fourier transform
of $f_{B,y}$ at $\pi_\lambda$ and at $\pi^w_{-\lambda}$.

So the function $\varphi_A(f,\pi)=-\operatorname{tr}(R^{-1}R'\,I(\pi,f))$
on the unitary dual of the cover of $A$, built from the normalized
operators $R$, has a variation which is the Fourier transform of the
variation of $\gamma\mapsto J_A(\gamma,f)$. Its inverse Fourier transform
$\Phi_A(\gamma,f)$ then varies exactly as $J_A(\gamma,f)$ does, and
\[
I_A(\gamma,f)=J_A(\gamma,f)-\Phi_A(\gamma,f)
\]
is invariant (Theorem~\ref{thm:IA}). The two functions live on opposite
sides of Fourier inversion: $\varphi_A(f,\cdot)$ on characters $\pi$ of
the cover of $A(\Q_S)$, $\Phi_A(\cdot,f)$ on its elements. Globally the
corrections on the two sides are matched by Poisson summation on the
cover of $A$, which is the rank-$0$ trace formula of
Section~\ref{sec:rank0}, so the identity is preserved. The normalized
operators $R_v$ are used, rather than $M$, because they are holomorphic
and unitary on $i\R$ and trivial on the unramified vectors, so that
$\varphi_A$ is given by finitely many local terms. The scalar that
remains, $m$ in $M=m\cdot R$, contributes $\frac{m'}m\operatorname{tr}
I(\pi,f)$, which is invariant; this is the term that survives on the
spectral side.

\paragraph{Results.} The main points of the section are these.
\begin{itemize}
\item The weight in the variance formulas is $u'_B(k,y)=-H(ky)$
  (Lemma~\ref{lem:uprime}).
\item The normalization of the intertwining operators is a theorem of
  Li~\cite{LiII}, and no trace Paley--Wiener theorem is needed for this
  group.
\item The invariant spectral terms attached to $\chi\in\mathcal X_A$ do
  not vanish. The logarithmic derivative $m'/m$ of the global normalizing
  factor survives, and $m$ is identified up to finitely many local
  factors: $m'/m=2\frac{L'}{L}(2\lambda,\chi^2)-2\frac{L'}{L}(2\lambda+1,
  \chi^2)$ plus the terms of the places in $S_{\mathrm{bad}}(\mu)$
  (Corollary~\ref{cor:global-m}).
\item The residual terms $\operatorname{tr}(R(f)|_{V_b})$ and the ramified
  terms of Theorem~\ref{thm:Jchi}(c) are already invariant. The ramified
  term is evaluated: $M(w,0)=-1$ on $\mathcal H_B(\chi_{\psi_b})$, so it
  equals $-\frac14\operatorname{tr}I(\chi_{\psi_b,0},f)$
  (Proposition~\ref{prop:Mw0}).
\item The function $\varphi_A(f)$ that carries the correction is not the
  Fourier transform of a compactly supported function on $\AS$
  (Example~\ref{ex:Phi-real}).
\end{itemize}

\paragraph{Conventions.} $S$ is a finite set of places of $\Q$ containing
$\infty$, $2$ and the places below which $\psi$ is ramified, $S_k$ is the
set of places of $k$ above $S$, and $\GS$, $\mathcal H(\GS)$,
$f=f_S\otimes e^S$ and $\gamma_S$ are as in Section~\ref{sec:local-dist}.
Let $K_S=\prod_{v\in S_k}K_v$, with preimage $\KS$ in $\GS$, and let
$\AS$ be the preimage of $A(\Q_S)=\prod_{v\in S_k}A(k_v)$, an abelian
group (it is the group $\overline{A(\A_S)}$ of Section~\ref{sec:IA}). Put
$H_S(x)=\sum_{v\in S_k}H_v(x_v)$ for $x\in\GS$; on the image of $\GS$ in
$\GA$ it is the height $H$ of Lemma~\ref{lem:NAK-global}. Let
$\Pi_u(\AS)$ be the group of genuine unitary characters of $\AS$. It is
a principal homogeneous space under the unitary dual of $A(\Q_S)$
(Lemma~\ref{lem:genuine-chars}), and $d\pi$ is the measure on it that
corresponds to the Plancherel measure dual to the Haar measure $da$ on
$A(\Q_S)$. For an anti-genuine integrable function $h$ on $\AS$ put
\[
\hat h(\pi)=\int_{\AS}h(\tilde a)\pi(\tilde a)\,d\tilde a,\qquad\text{so
that}\qquad h(\tilde a)=\int_{\Pi_u(\AS)}\hat h(\pi)\,\pi(\tilde a)^{-1}\,d\pi
\]
when $\hat h$ is integrable and $h$ continuous. For $\pi\in\Pi_u(\AS)$ and
$\lambda\in\C$ put $\pi_\lambda=\pi\,e^{\lambda H_S}$, and
$\pi^w(\tilde a)=\pi(w^{-1}\tilde aw)$, where $w$ stands for
$\prod_{v\in S_k}(w,1)_v$; by Lemma~\ref{lem:w-action}(a),
$(\pi_\lambda)^w=(\pi^w)_{-\lambda}$. Let $V(\pi)$ be the space of smooth
functions $\varphi$ on $\KS$ with $\varphi(\tilde m\tilde k)=\pi(\tilde
m)\varphi(\tilde k)$ for $\tilde m\in\AS\cap\KS$, and let $I(\pi_\lambda)$
be the induced representation of $\GS$ on $V(\pi)$:
$(I(\pi_\lambda,y)\varphi)(\tilde k)=\varphi_\lambda(\tilde ky)$, where
$\varphi_\lambda(n\tilde a\tilde k)=\pi(\tilde a)e^{(\lambda+1)H_S(\tilde
a)}\varphi(\tilde k)$. These are the local analogs of
$\mathcal H_B(\mu)$ and $I(\mu_\lambda)$ in
Section~\ref{sec:eisenstein}: for $\mu\in\Pi_{\mathrm{aut}}(X_A)$ with
$\mu_v$ trivial on $A(\OO_v)^*$ for all $v\notin S_k$, the restriction
$\mu_S$ of $\mu$ to $\AS$ lies in $\Pi_u(\AS)$, and
$\operatorname{tr}I(\mu_\lambda,f)=\operatorname{tr}I(\mu_{S,\lambda},f_S)$.

Invariance refers to conjugation by all $y\in\GS$. At an archimedean
place $f^y$ need not be $K$-finite when $f$ is, so in
Sections~\ref{sec:noninv} and~\ref{sec:IAgamma} test functions are taken
in $C_c^\infty(\GS)_-$, the space of smooth, compactly supported,
anti-genuine functions, which is stable under $f\mapsto f^y$ and to which
the distributions $J_G(\gamma,\cdot)$, $J_A(\gamma,\cdot)$ and
$\operatorname{tr}I(\pi,\cdot)$ extend by the same formulas.

\subsection{Non-invariance}\label{sec:noninv}

This subsection computes $J(f^y)-J(f)$ for each of the distributions in
the trace formula. In every case the difference is a rank-$0$
distribution on the cover of $A$, evaluated at the function $f_{B,y}$
defined below. The geometric and the spectral differences are the two
sides of the trace formula for $\overline{A(\A)^1}$
(Theorem~\ref{thm:tf-rank0}), which is why the correction terms of
Section~\ref{sec:itf} cancel in the sum.

As in~\cite[\S\S2--3]{A6}, let
\begin{equation}\label{eq:Gammaprime}
\Gamma'_B(X,Y)=\tau_B(X)-\hat\tau_B(X-Y),\qquad X,Y\in\mathfrak a\cong\R,
\end{equation}
the case $Q=B$ of $\Gamma'_Q(X,Y)=\sum_{R\supset Q}(-1)^{\dim(A_R/A_G)}
\tau_Q^R(X)\hat\tau_R(X-Y)$, in the notation of~\eqref{eq:tau}; for fixed
$Y$ it is compactly supported in $X$. For $\tilde k\in\KS$, $y\in\GS$ and
$f\in C_c^\infty(\GS)_-$ put
\begin{gather}
u'_B(\tilde k,y)=\int_{\mathfrak a}\Gamma'_B\bigl(X,-H_S(\tilde ky)\bigr)
\,dX,\label{eq:uprime}\\
f_{B,y}(\tilde a)=\delta_B(\tilde a)^{1/2}\int_{\KS}\int_{N(\Q_S)}
f(\tilde k^{-1}\tilde an\tilde k)\,u'_B(\tilde k,y)\,dn\,d\tilde k,\qquad
\tilde a\in\AS,\label{eq:fBy}
\end{gather}
and likewise for $y\in\GA$, $f\in\mathcal H(\GA)$ and $\tilde a\in
\overline{A(\A)}$, with $\overline{K_0}$, $N(\A)$ and $H$. Here $dX$ is
Lebesgue measure, as in Remark~\ref{rem:alpha}.

\begin{lemma}\label{lem:uprime}
\begin{enumerate}[label=(\alph*)]
\item $u'_B(\tilde k,y)=-H_S(\tilde ky)$.
\item $f_{B,y}$ is smooth, compactly supported and anti-genuine on $\AS$,
  and $f_{B,1}=0$.
\item For $\pi\in\Pi_u(\AS)$ and $\lambda\in\C$ let $U_y$ be the operator
  of multiplication by $\tilde k\mapsto u'_B(\tilde k,y)$ on $V(\pi)$.
  Then $\operatorname{tr}\bigl(U_y\,I(\pi_\lambda,f)\bigr)
  =\widehat{f_{B,y}}(\pi_\lambda)$.
\end{enumerate}
\end{lemma}

\begin{proof}
(a) $\int_\R\bigl(\tau_B(X)-\tau_B(X-Y)\bigr)\,dX=Y$ for every real $Y$:
the integrand is the characteristic function of $(0,Y]$ if $Y>0$ and
minus that of $(Y,0]$ if $Y<0$. (b) $u'_B(\cdot,y)$ is smooth and left
invariant under $\AS\cap\KS$, and $u'_B(\tilde k,1)=-H_S(\tilde k)=0$.
(c) As in the proof of Proposition~\ref{prop:Tcoef}, $I(\pi_\lambda,f)$
has the kernel $\int_{\AS}\int_{N(\Q_S)}f(\tilde k^{-1}\tilde
an\tilde k')\pi(\tilde a)e^{(\lambda+1)H_S(\tilde a)}\,dn\,d\tilde a$ on
$\KS$; multiply by $u'_B(\tilde k,y)$ and integrate over the diagonal,
using $\delta_B^{1/2}=e^{H_S}$.
\end{proof}

The sign matters: $u'_B(\cdot,y)$ is not of constant sign
(Example~\ref{ex:uprime}), and with $|H|$ in place of $-H$ the formulas of
Lemmas~\ref{lem:variance-global} and~\ref{lem:variance-local} fail.

\begin{example}\label{ex:uprime}
Let $v$ be a real place, $k_\theta=\begin{pmatrix}\cos\theta&\sin\theta\\
-\sin\theta&\cos\theta\end{pmatrix}$ and $y=a(s)$ with $s>0$. Since
$H_v(g)=-\log\|e_2g\|$ (Lemma~\ref{lem:height-ineq}),
\[
u'_B(k_\theta,a(s))=\tfrac12\log\bigl(s^2\sin^2\theta+s^{-2}\cos^2\theta
\bigr),
\]
which is $-\log s$ at $\theta=0$ and $+\log s$ at $\theta=\frac\pi2$.
\end{example}

For a class $\mathcal O\in\mathcal C$, a class $\chi\in\mathcal X$ and an
anti-genuine function $h$ on $\overline{A(\A)}$ whose restriction to
$\overline{A(\A)^1}$ is smooth and compactly supported, put
\begin{equation}\label{eq:JAo}
J^A_{\mathcal O}(h)=c_k\sum_{\gamma\in\mathcal O\cap A(\Q)}h(\gamma),\qquad
J^A_\chi(h)=\sum_{\mu\in\chi}\int_{\overline{A(\A)^1}}h(\tilde a)\mu(\tilde
a)\,d\tilde a,
\end{equation}
with $J^A_\chi=0$ for $\chi\in\mathcal X_G$. These are the
distributions~\eqref{eq:IA}, summed over $\mathcal O\cap A(\Q)$ and over
$\chi$, and Theorem~\ref{thm:tf-rank0} says that
$\sum_{\mathcal O}J^A_{\mathcal O}(h)=\sum_\chi J^A_\chi(h)$. By Fourier
inversion on $\mathfrak a\cong\R$, $\int_{\overline{A(\A)^1}}h\mu=\frac1{2\pi}
\int_\R\hat h(\mu_{it})\,dt$ when $\hat h(\mu_{it})=\int_{\overline{A(\A)}}
h\,\mu\,e^{itH}$ is integrable in $t$.

\begin{lemma}\label{lem:variance-global}
Let $f\in\mathcal H(\GA)$ and $y\in\GA$ with $f^y\in\mathcal H(\GA)$. Then,
for all $\mathcal O\in\mathcal C$ and $\chi\in\mathcal X$,
\[
J_{\mathcal O}(f^y)=J_{\mathcal O}(f)+J^A_{\mathcal O}(f_{B,y}),\qquad
J_\chi(f^y)=J_\chi(f)+J^A_\chi(f_{B,y}).
\]
In particular $J_{\mathcal O}$ is invariant when $\mathcal O$ is elliptic,
and $J_\chi$ is invariant when $\chi\in\mathcal X_G$.
\end{lemma}

\begin{proof}
This is the formula of~\cite[\S3]{A6} for this group; the proof is
short. Since
$f^y(x^{-1}\gamma nx)=f\bigl((xy^{-1})^{-1}\gamma n(xy^{-1})\bigr)$, the
kernels satisfy $K_{P,\mathcal O}^{f^y}(x,x)=K_{P,\mathcal O}^f(xy^{-1},
xy^{-1})$. Substituting $x\mapsto xy$ in $J^T_{\mathcal O}(f^y)=\int
k^T_{\mathcal O}(x,f^y)\,dx$ (Theorem~\ref{thm:coarse}) replaces
$\hat\tau_B(H(\delta x)-T)$ in~\eqref{eq:kT} by $\hat\tau_B(H(\delta
xy)-T)$. Write $\delta x=n\tilde a\tilde k$; then $H(\delta xy)=H(\delta
x)+H(\tilde ky)$, and by~\eqref{eq:Gammaprime},
\[
\hat\tau_B\bigl(H(\delta x)+H(\tilde ky)-T\bigr)=\hat\tau_B\bigl(H(\delta
x)-T\bigr)-\Gamma'_B\bigl(H(\delta x)-T,-H(\tilde ky)\bigr).
\]
So $J^T_{\mathcal O}(f^y)-J^T_{\mathcal O}(f)=\int_{B(\Q)\backslash\GA}
K_{B,\mathcal O}(x,x)\,\Gamma'_B(H(x)-T,-H(\tilde ky))\,dx$. In the
coordinates $x=n\tilde a\tilde k$, with $dx=\delta_B(\tilde a)^{-1}dn\,
d\tilde a\,d\tilde k$,
\[
K_{B,\mathcal O}(x,x)=\delta_B(\tilde a)\sum_{\gamma\in\mathcal O\cap A(\Q)}
\int_{N(\A)}f(\tilde k^{-1}\gamma n\tilde k)\,dn,
\]
and $N(\Q)\backslash N(\A)$ has volume $1$. The integral over
$A(\Q)\backslash\overline{A(\A)}$ of $\Gamma'_B(H(\tilde a)-T,-H(\tilde
ky))$ is $c_k\,u'_B(\tilde k,y)$, for every $T$. Since $\delta_B(\gamma)=1$
for $\gamma\in A(\Q)$, the difference is $c_k\sum_\gamma f_{B,y}(\gamma)$.
All integrals converge absolutely: the sum over $\gamma$ is finite and
$\Gamma'_B(\cdot,Y)$ is compactly supported.

For $\chi$ the argument is the same, with~\eqref{eq:kchi} and
$K_{B,\chi}$. By~\eqref{eq:KBchi}, and because an orthonormal basis
for~\eqref{eq:HBmu} satisfies $c_k\int_{\overline{K_0}}|\varphi|^2=1$, the
difference is
\[
\sum_{\mu\in\chi}\frac1{2\pi}\int_\R\operatorname{tr}\bigl(U_y\,
I(\mu_{it},f)\bigr)\,dt=\sum_{\mu\in\chi}\frac1{2\pi}\int_\R
\widehat{f_{B,y}}(\mu_{it})\,dt=J^A_\chi(f_{B,y})
\]
by Lemma~\ref{lem:uprime}(c); the interchange of integrals is justified
as in~\cite[\S3]{A6}, by Theorem~\ref{thm:Jchi-poly}.
\end{proof}

Written as a sum over all parabolic subgroups $Q\supset A$, as
in~\cite{A9,A10}, the formula is $J(f^y)=\sum_Q|W_0^{M_Q}||W_0|^{-1}
J^{M_Q}(f_{Q,y})$. Here $Q$ runs over $G$, $B$ and $\bar B$; the last
two contribute $\frac12J^A(f_{B,y})$ each, because
$J^A(f_{\bar B,y})=J^A(f_{B,y})$.

The local distributions vary in the same way. Recall $J_G$, $J_A$ and the
weight $v_A$ from~\eqref{eq:JG}--\eqref{eq:JA}, and $J_A(z,f)$ from
Lemma~\ref{lem:JA-central}.

\begin{lemma}\label{lem:variance-local}
Let $f\in C_c^\infty(\GS)_-$ and $x,y\in\GS$.
\begin{enumerate}[label=(\alph*)]
\item If $x=n\tilde a\tilde k$ and $wx=n'\tilde a'\tilde k'$ are Iwasawa
  decompositions, then
  $v_A(xy)=v_A(x)-H_S(\tilde ky)-H_S(\tilde k'y)$.
\item If $\gamma\in\AS$ is regular (that is, $\gamma$ lies over $a(t)$
  with $t_v^2\ne1$ for all $v\in S_k$), then $J_G(\gamma,f^y)=J_G(\gamma,f)$
  and
  \[
  J_A(\gamma,f^y)=J_A(\gamma,f)+f_{B,y}(\gamma)+f_{B,y}(w\gamma w^{-1}).
  \]
\item For $z\in\{\pm1\}$, $J_A(z,f^y)=J_A(z,f)+2f_{B,y}(z)$.
\end{enumerate}
\end{lemma}

\begin{proof}
(a) $H_S(xy)=H_S(x)+H_S(\tilde ky)$ and $H_S(wxy)=H_S(wx)+H_S(\tilde
k'y)$; subtract from $v_A(x)=-H_S(x)-H_S(wx)$. (b) Substituting $x\mapsto
xy$ in~\eqref{eq:JA} gives $J_A(\gamma,f^y)=|D(\gamma)|^{1/2}\int
f(x^{-1}\gamma x)v_A(xy)\,dx$, and likewise for $J_G$ without the weight.
By~(a) the difference from $J_A(\gamma,f)$ is the sum of two terms. In the
first, with $x=n\tilde k$ and $n^{-1}\gamma n=\gamma n'$, where
$dn=|1-t^{-2}|^{-1}dn'$ and $|D(\gamma)|^{1/2}=\delta_B(\gamma)^{1/2}
|1-t^{-2}|$, we get $f_{B,y}(\gamma)$. In the second, substitute
$x\mapsto w^{-1}x$; this replaces $\gamma$ by $w\gamma w^{-1}$ and
$\tilde k'$ by the $\KS$-component of $x$, and
$D(w\gamma w^{-1})=D(\gamma)$. (c) Apply~(b) to $a_t$ in
Lemma~\ref{lem:JA-central}(b): $J_G(a_t,\cdot)$ and $r(t)$ do not change,
and $f_{B,y}(a_t)+f_{B,y}(wa_tw^{-1})\to2f_{B,y}(z)$ as $t\to1$, because
$f_{B,y}$ is continuous and $z$ is central.
\end{proof}

Parts~(b) and~(c) are the case $M=A$ of Arthur's formula
$J_M(\gamma,f^y)=\sum_{Q\in\mathcal F(M)}J_M^{M_Q}(\gamma,f_{Q,y})$
\cite{A6,A9}, with $\mathcal F(A)=\{G,B,\bar B\}$ and
$f_{\bar B,y}(\gamma)=f_{B,y}(w\gamma w^{-1})$.

On the spectral side everything follows from one identity. It holds for
the global operator $M(w,\lambda)$, for the local operators and for their
normalized versions alike, so we state it for any intertwining family.

\begin{lemma}\label{lem:conj-identity}
Let $\pi\in\Pi_u(\AS)$, and let $\mathcal M(\lambda)\colon V(\pi)\to
V(\pi^w)$ be a meromorphic family of invertible operators with
$\mathcal M(\lambda)I(\pi_\lambda,y)=I(\pi^w_{-\lambda},y)\mathcal
M(\lambda)$ for all $y\in\GS$. For $y\in\GS$ let $H_y$ be multiplication
by $\tilde k\mapsto H_S(\tilde ky)$, on $V(\pi)$ and on $V(\pi^w)$. Then,
with $\mathcal M'=d\mathcal M/d\lambda$,
\begin{equation}\label{eq:conj-identity}
I(\pi_\lambda,y)\,\mathcal M(\lambda)^{-1}\mathcal M'(\lambda)\,
I(\pi_\lambda,y)^{-1}=\mathcal M(\lambda)^{-1}\mathcal M'(\lambda)+H_y
+\mathcal M(\lambda)^{-1}H_y\,\mathcal M(\lambda),
\end{equation}
and for $f\in C_c^\infty(\GS)_-$,
\begin{equation}\label{eq:trace-variance}
-\operatorname{tr}\bigl(\mathcal M^{-1}\mathcal M'\,I(\pi_\lambda,f^y)\bigr)
=-\operatorname{tr}\bigl(\mathcal M^{-1}\mathcal M'\,I(\pi_\lambda,f)\bigr)
+\widehat{f_{B,y}}(\pi_\lambda)+\widehat{f_{B,y}}(\pi^w_{-\lambda}).
\end{equation}
The same holds over $\A$, for $\mathcal H_B(\mu)$, $I(\mu_\lambda)$ and
$\GA$.
\end{lemma}

\begin{proof}
In the compact picture, $(I(\pi_\lambda,y)\varphi)(\tilde k)=\pi(\tilde
a)e^{(\lambda+1)H_S(\tilde ky)}\varphi(\tilde k_1)$ for $\tilde
ky=n\tilde a\tilde k_1$, so $\frac{d}{d\lambda}I(\pi_\lambda,y)=H_y
I(\pi_\lambda,y)$ and $\frac{d}{d\lambda}I(\pi^w_{-\lambda},y)=-H_y
I(\pi^w_{-\lambda},y)$. Differentiating the intertwining relation,
\[
\mathcal M'I(\pi_\lambda,y)+\mathcal MH_yI(\pi_\lambda,y)
=-H_yI(\pi^w_{-\lambda},y)\mathcal M+I(\pi^w_{-\lambda},y)\mathcal M'.
\]
Multiply by $\mathcal M^{-1}$ on the left and by $I(\pi_\lambda,y)^{-1}$
on the right, and use $\mathcal M^{-1}I(\pi^w_{-\lambda},y)=I(\pi_\lambda,
y)\mathcal M^{-1}$; this is~\eqref{eq:conj-identity}. Since
$I(\pi_\lambda,f^y)=I(\pi_\lambda,y)^{-1}I(\pi_\lambda,f)I(\pi_\lambda,y)$,
the left side of~\eqref{eq:trace-variance} is $-\operatorname{tr}\bigl(
I(\pi_\lambda,y)\mathcal M^{-1}\mathcal M'I(\pi_\lambda,y)^{-1}
I(\pi_\lambda,f)\bigr)$. Insert~\eqref{eq:conj-identity}; the last term
contributes $-\operatorname{tr}(H_y\,\mathcal MI(\pi_\lambda,f)\mathcal
M^{-1})=-\operatorname{tr}(H_yI(\pi^w_{-\lambda},f))$, and $-H_y=U_y$, so
Lemma~\ref{lem:uprime}(c) applies to both terms.
\end{proof}

With $\mathcal M=M(w,\lambda)$, formula~\eqref{eq:trace-variance},
Theorem~\ref{thm:Jchi} and the invariance of $\operatorname{tr}(R(f)|_{V_b})$
and of $\operatorname{tr}(M(w,0)I(\mu_0,f))$ give a second proof of the
formula for $J_\chi(f^y)$ in Lemma~\ref{lem:variance-global}. (The
operator $M(w,0)$ commutes with $I(\mu_0,y)$ when $\mu^w=\mu$, so
$f\mapsto\operatorname{tr}(M(w,0)I(\mu_0,f))$ is invariant.)

\subsection{Normalized intertwining operators}\label{sec:normalized}

The correction terms are built from the logarithmic derivative of the
intertwining operator. The global operator $M(w,\lambda)$ is not a
product of local operators that are holomorphic on $i\R$: the local
operators have poles there, and the Euler product of their unramified
eigenvalues does not converge. One divides each local operator by a
scalar normalizing factor; the product of these factors is the global
function $m(\mu,\lambda)$, and what remains is a finite product of
normalized local operators, which are unitary and holomorphic on $i\R$.

Let $v$ be a place of $k$, $\pi_v$ a genuine unitary character of
$\overline{A(k_v)}$, and
\[
\bigl(M_v(\pi_v,\lambda)\varphi\bigr)_{-\lambda}(x)=\int_{k_v}
\varphi_\lambda\bigl((w,1)^{-1}(n(u),1)x\bigr)\,du,\qquad
M_v(\pi_v,\lambda)\colon V(\pi_v)\to V(\pi_v^w),
\]
as in Lemma~\ref{lem:intertwining}; the integral converges for
$\mathrm{Re}\,\lambda>0$ and continues meromorphically.

\begin{theorem}[Arthur~{\cite[Thm.~2.1]{A11}}; Li~{\cite[\S3]{LiII}} for
covers]\label{thm:normalization}
There are meromorphic functions $r_v(\pi_v,\lambda)$ of $\lambda\in\C$,
with $r_v(\pi_{v,\nu},\lambda)=r_v(\pi_v,\lambda+\nu)$ for $\nu\in i\R$,
such that the operators $R_v(\pi_v,\lambda)=r_v(\pi_v,\lambda)^{-1}
M_v(\pi_v,\lambda)$ have the following properties.
\begin{enumerate}[label=(N\arabic*)]
\item\label{N1} $R_v(\pi_v,\lambda)I(\pi_{v,\lambda},x)=I(\pi^w_{v,-\lambda},x)
  R_v(\pi_v,\lambda)$ for $x\in\overline{SL(2,k_v)}$.
\item\label{N2} $R_v(\pi_v^w,-\lambda)R_v(\pi_v,\lambda)=1$, and for
  $\lambda\in i\R$ the operator $R_v(\pi_v,\lambda)$ is holomorphic and
  unitary.
\item\label{N3} On each $\overline{K_v}$-isotypic subspace,
  $R_v(\pi_v,\lambda)$ is a rational function of $q_v^{-\lambda}$ if $v$
  is finite, and a rational function of $\lambda$ if $v$ is archimedean.
  In the archimedean case the norms of $R_v^{-1}\frac{d}{d\lambda}R_v$ on
  the $\overline{K_v}$-types, for $\lambda\in i\R$, are bounded by a
  polynomial in the highest weight of the type and in $|\lambda|$ (for
  real $v$ this is explicit in~\eqref{eq:rho-m}).
\item\label{N4} If $v$ is finite and non-dyadic, $\psi_v$ is unramified
  and $\pi_v$ is trivial on $A(\OO_v)^*$, then
  $R_v(\pi_v,\lambda)\varphi^0_v=\varphi^{0,w}_v$ for all $\lambda$.
\end{enumerate}
\end{theorem}

These are Arthur's conditions $(R_1)$--$(R_8)$ of~\cite[\S2]{A11} for the
pair $(B,\bar B)$, written with $w$: $M_v(\pi_v,\lambda)$ is the
composition of the operator $J_{\bar B|B}(\pi_{v,\lambda})$ of~\cite{A11}
with left translation by $w$, which does not depend on $\lambda$.
At the real places everything is explicit for the double cover
(Lemma~\ref{lem:explicit-r}(b),(c)); for the universal cover of
$SL(2,\R)$ see Sally~\cite{Sa}. The factors $r_v$ are not
unique. Three statements about them should be kept apart.
\begin{enumerate}[label=(\roman*)]
\item \emph{Existence.} At every place there are factors $r_v$ with
  \ref{N1}--\ref{N4}. This is Theorem~\ref{thm:normalization}, and it is
  all that the construction of the invariant distributions uses.
\item \emph{An explicit formula at a given place.} We have one at the
  finite places covered by~\ref{N4}, at the real places and at the
  complex places (Lemma~\ref{lem:explicit-r} and the remark after
  Example~\ref{ex:real-R}). We do not have one at the places of
  $S_{\mathrm{bad}}(\mu)$ defined below (Remark~\ref{rem:bad-places}).
\item \emph{Compatibility with the global operator.} For any choice of
  local factors, the product $m(\mu,\lambda)$ in~\eqref{eq:global-m} is
  defined, $M(w,\lambda)=m\cdot R_S$, and $m$ has the properties of
  Corollary~\ref{cor:global-m}(b). So no compatibility condition has to be
  imposed on the local choices; but $m$, $m'/m$ and the distributions
  $I_A(\gamma,f)$ depend on them, and $m$ is a known function only up to
  the unknown factors at $S_{\mathrm{bad}}(\mu)$.
\end{enumerate}
At most places there is a natural choice:

\begin{lemma}\label{lem:explicit-r}
Write $\pi_v=\chi_{\psi,v}\cdot(\chi_v\circ p)$.
\begin{enumerate}[label=(\alph*)]
\item Let $v$, $\psi_v$ and $\pi_v$ be as in~\ref{N4}. Then
  \[
  r_v(\pi_v,\lambda)=\frac{L_v(2\lambda,\chi_v^2)}{L_v(2\lambda+1,\chi_v^2)}
  \]
  satisfies~\ref{N1}--\ref{N4}.
\item Let $v$ be real and $\chi_v=\mathrm{sgn}^\epsilon$,
  $\epsilon\in\{0,1\}$. Let $\theta\mapsto\tilde k_\theta$ be the
  one-parameter subgroup of $\overline{K_v}$ above $k_\theta$
  (Example~\ref{ex:uprime}). The $\overline{K_v}$-types of $I(\pi_v)$ are
  the characters $\tilde k_\theta\mapsto e^{im\theta}$, $m\in m_0+2\Z$,
  each with multiplicity one, where $m_0\in\{\frac12,-\frac12\}$ is
  determined by $e^{i\pi m_0}=\pi_v(a(-1),1)$. On the line of weight $m$,
  with $du$ Lebesgue measure,
  \begin{equation}\label{eq:real-M}
  M_v(\pi_v,\lambda)=e^{i\pi m/2}\,\frac{\pi\,2^{1-\lambda}\,\Gamma(\lambda)}
  {\Gamma\bigl(\frac{\lambda+1+m}2\bigr)\Gamma\bigl(\frac{\lambda+1-m}2\bigr)} .
  \end{equation}
  With $r_v(\pi_v,\lambda)=\sqrt2\,e^{i\pi m_0/2}\,L_v(2\lambda)/L_v(2\lambda+1)$,
  where $L_v(s)=\pi^{-s/2}\Gamma(s/2)$, the operator $R_v(\pi_v,\lambda)$
  is the identity on the weight $m_0$, and on the weight $m_0+2j$ it is
  \begin{equation}\label{eq:real-R}
  \prod_{i=0}^{j-1}\frac{a_i-\lambda}{a_i+\lambda}\quad(j>0,\ a_i=1+m_0+2i),
  \qquad
  \prod_{i=0}^{|j|-1}\frac{b_i-\lambda}{b_i+\lambda}\quad(j<0,\ b_i=1-m_0+2i).
  \end{equation}
  These operators satisfy~\ref{N1}--\ref{N3}.
\item Let $v$ be real. Every genuine unitary character of
  $\overline{A(k_v)}$ is of the form $\pi_{v,i\tau}=\pi_v\,e^{i\tau H_v}$
  with $\pi_v$ as in~(b) and a unique $\tau\in\R$; that is,
  $\chi_v=\mathrm{sgn}^\epsilon|\cdot|^{i\tau}$. Then
  $V(\pi_{v,i\tau})=V(\pi_v)$, $I(\pi_{v,i\tau},\lambda)=I(\pi_v,\lambda+
  i\tau)$ and $M_v(\pi_{v,i\tau},\lambda)=M_v(\pi_v,\lambda+i\tau)$. With
  $r_v(\pi_{v,i\tau},\lambda)=r_v(\pi_v,\lambda+i\tau)$, the normalized
  operator $R_v(\pi_{v,i\tau},\lambda)$ is given by~\eqref{eq:real-R} with
  $\lambda+i\tau$ in place of $\lambda$, and satisfies~\ref{N1}--\ref{N3}.
  In terms of $\chi_v$,
  \[
  r_v(\pi_{v,i\tau},\lambda)=\sqrt2\,e^{i\pi m_0/2}\,
  \frac{L_v(2\lambda,\chi_v^2)}{L_v(2\lambda+1,\chi_v^2)},\qquad
  L_v(s,|\cdot|^{2i\tau})=L_v(s+2i\tau).
  \]
\end{enumerate}
\end{lemma}

\begin{proof}
(a) \ref{N1} holds for $M_v$, and~\ref{N4} is
Lemma~\ref{lem:intertwining}(b). The induced representation
$I(\pi_{v,\lambda})$ is irreducible for $\lambda$ outside a discrete
set, so $M_v(\pi_v^w,-\lambda)M_v(\pi_v,\lambda)$, which commutes with
it, is a scalar; on $\varphi_v^0$ it is $r_v(\pi_v^w,-\lambda)
r_v(\pi_v,\lambda)$ by Lemma~\ref{lem:intertwining}(b). So
$R_v(\pi^w_v,-\lambda)R_v(\pi_v,\lambda)=1$. For $\lambda\in i\R$ the
adjoint of $M_v(\pi_v,\lambda)$ is $M_v(\pi_v^w,-\lambda)$~\cite[\S1]{A11}. (This
uses that $(w,1)_v^2$ acts trivially on $I(\pi_{v,\lambda})$. Here it
does: $(-1,-1)_v=1$ at a non-dyadic finite place, so $(w,1)_v^2=(-1,1)$,
a central element lying in $A(\OO_v)^*$, and it acts by the scalar
$\pi_v(-1,1)=1$. At a real place $(-1,1)$ has order $4$ and acts by $\pm
i$ on a genuine representation, by Lemma~\ref{lem:local}(e); part~(b)
does not use the adjoint formula.) Moreover
$\overline{r_v(\pi_v,\lambda)}=r_v(\pi_v^w,-\lambda)$ for $\lambda\in i\R$;
hence $R_v$ is unitary on $i\R$ wherever it is defined, so it is bounded
there and has no poles. Rationality is~\cite{LiII}.

(b) In the model of Section~\ref{sec:local}, $\tilde k_\theta=(k_\theta,1)$
for $0<\theta<\pi$: for $\theta_1,\theta_2,\theta_1+\theta_2\in(0,\pi)$
the three values of $x(\cdot)$ in~\eqref{eq:beta} are $-\sin$ of the
angles, all negative, so $\beta(k_{\theta_1},k_{\theta_2})=1$. Then
$\tilde k_\pi=(k_{\pi/2},1)^2=(-1,1)=(a(-1),1)$, which gives the
$\overline{K_v}$-types, since $\overline{A(k_v)}\cap\overline{K_v}$ is
generated by $\tilde k_\pi$. Also $(w,1)^{-1}=\tilde k_{\pi/2}$, and
$\tilde k_{\pi/2}(n(u),1)=(n',1)(a(s),1)\tilde k_{\theta(u)}$ with
$s=(1+u^2)^{-1/2}>0$ and $\theta(u)=\frac\pi2+\arctan u$, by
continuity in $u$. So on a vector of weight $m$ the
integral is $e^{i\pi m/2}\int_\R(1+u^2)^{-(\lambda+1)/2}e^{im\arctan u}\,du$,
and $\int_{-\pi/2}^{\pi/2}\cos^{\lambda-1}\vartheta\,e^{im\vartheta}\,
d\vartheta$ is the classical integral in~\eqref{eq:real-M}. By the
duplication formula, $\Gamma(\frac{\lambda}2+\frac34)\Gamma(\frac\lambda2+
\frac14)=2^{\frac12-\lambda}\sqrt\pi\,\Gamma(\lambda+\frac12)$,
so~\eqref{eq:real-M} at $m=\pm\frac12$ is $\sqrt2\,e^{i\pi m/2}\sqrt\pi\,
\Gamma(\lambda)/\Gamma(\lambda+\frac12)$. The quotient of~\eqref{eq:real-M}
at $m+2$ and at $m$ is $-(\lambda-1-m)/(\lambda+1+m)$,
which gives~\eqref{eq:real-R}. Each factor in~\eqref{eq:real-R} has
absolute value $1$ on $i\R$ and no pole there, because $a_i$ and $b_i$
are nonzero real numbers; \ref{N2} and~\ref{N3} follow, and~\ref{N1}
holds for $M_v$.

(c) The unitary characters of $\R^\times$ are the
$\mathrm{sgn}^\epsilon|\cdot|^{i\tau}$, and $H_v(a(t),\zeta)=\log|t|$. The
space $V(\pi)$ depends only on the restriction of $\pi$ to
$\overline{A(k_v)}\cap\overline{K_v}$, on which $e^{i\tau H_v}=1$, and
$\varphi_\lambda$ formed with $\pi_{v,i\tau}$ is $\varphi_{\lambda+i\tau}$
formed with $\pi_v$; this gives the first three statements. The factors
in~\eqref{eq:real-R} with $\lambda+i\tau$ in place of $\lambda$ still have
absolute value $1$ and no poles for $\lambda\in i\R$. The last formula is
$L_v(s,|\cdot|^{2i\tau})=\pi^{-(s+2i\tau)/2}\Gamma((s+2i\tau)/2)$.
\end{proof}

\begin{example}\label{ex:real-R}
In Lemma~\ref{lem:explicit-r}(b) take $m_0=\frac12$, which is the case
$\pi_v=\chi_{\psi,v}$ for $\psi_v(x)=e^{-2\pi ix}$
(Example~\ref{ex:weil-Q}). The weights are $\frac12,\frac52,\frac92,\dots$
and $-\frac32,-\frac72,\dots$, and
\[
R_v(\lambda)\big|_{5/2}=\frac{\frac32-\lambda}{\frac32+\lambda},\qquad
R_v(\lambda)\big|_{-3/2}=\frac{\frac12-\lambda}{\frac12+\lambda},\qquad
R_v^{-1}R_v'(\lambda)\big|_{5/2}=-\frac{3}{\frac94-\lambda^2}.
\]
$R_v(\lambda)$ has a zero at $\lambda=\frac12$ on every weight below
$\frac12$ and is $1$ on the weight $\frac12$: this is the reducibility of
$I(\chi_{\psi,v,1/2})$, whose quotient with lowest weight $\frac12$ is
the even Weil representation of Theorem~\ref{thm:residual}. In
general, by~\eqref{eq:real-R}, on the weight $m_0+2j$ and for $\lambda=it$,
\begin{equation}\label{eq:rho-m}
-R_v^{-1}R_v'(it)=\rho_j(t):=\sum_i\frac{2c_i}{c_i^2+t^2}\;\ge0,
\end{equation}
the sum over $c_i=a_i$, $0\le i<j$, if $j>0$, and over $c_i=b_i$,
$0\le i<|j|$, if $j<0$; and $\rho_0=0$. For $\pi_{v,i\tau}$ as in
Lemma~\ref{lem:explicit-r}(c) the value is $\rho_j(t+\tau)$. For $SL(2,\R)$ the weights are
integers and the same computation gives the familiar formula with
$L_v(\lambda)/L_v(\lambda+1)$.
\end{example}

At a complex place the cover splits and $R_v$ is the normalized operator
of $SL(2,\C)$~\cite{A11}; the factor
$L_v(\lambda,\chi_v)/L_v(\lambda+1,\chi_v)$ there equals $2L_v(2\lambda,
\chi_v^2)/L_v(2\lambda+1,\chi_v^2)$. So the factor of
Lemma~\ref{lem:explicit-r} is available at every place outside the finite
set
\[
S_{\mathrm{bad}}(\mu)=\{v\text{ dyadic}\}\cup\{v\mid\psi_v\text{ ramified}\}
\cup\{v\text{ finite}\mid\mu_v\text{ nontrivial on }A(\OO_v)^*\}.
\]

\paragraph{The global normalizing factor.} Let
$\mu=\chi_\psi\cdot(\chi\circ p)\in\Pi_{\mathrm{aut}}(X_A)$, and fix
factors $r_v(\mu_v,\lambda)$ as in Theorem~\ref{thm:normalization} for
$v\in S_{\mathrm{bad}}(\mu)$ and as in Lemma~\ref{lem:explicit-r} (and
the remark on complex places) for all other $v$. At a real place the
component $\mu_v$ is $\chi_{\psi,v}\cdot(\chi_v\circ p)$ with
$\chi_v=\mathrm{sgn}^\epsilon|\cdot|^{i\tau_v}$, and
Lemma~\ref{lem:explicit-r}(c) applies; for $k=\Q$ one has $\tau_\infty=0$,
because $\mu$ is trivial on $A_B(\R)^0$, but for other fields the
$\tau_v$ need not vanish. For a finite set
$S_1\supset S_{\mathrm{bad}}(\mu)$ of places of $k$ containing the
archimedean ones, put
\begin{equation}\label{eq:global-m}
m(\mu,\lambda)=m^{S_1}(\mu,\lambda)\prod_{v\in S_1}r_v(\mu_v,\lambda),
\qquad m^{S_1}(\mu,\lambda)=\frac{L^{S_1}(2\lambda,\chi^2)}
{L^{S_1}(2\lambda+1,\chi^2)}
\end{equation}
as in~\eqref{eq:mS}. By Lemma~\ref{lem:explicit-r}(a) this does not
depend on $S_1$.

\begin{corollary}\label{cor:global-m}
Let $f=f_S\otimes e^S$ and $\mu\in\Pi_{\mathrm{aut}}(X_A)$ with $\mu_v$
trivial on $A(\OO_v)^*$ for $v\notin S_k$, and let $R_S(\mu_S,\lambda)=
\bigotimes_{v\in S_k}R_v(\mu_v,\lambda)$.
\begin{enumerate}[label=(\alph*)]
\item On $\bigotimes_{v\in S_k}V(\mu_v)\otimes\bigotimes_{v\notin S_k}
  \varphi_v^0$, $M(w,\lambda)=m(\mu,\lambda)\,R_S(\mu_S,\lambda)\otimes1$,
  and for $\lambda\in i\R$
  \[
  \operatorname{tr}\bigl(M(w,\lambda)^{-1}M'(w,\lambda)I(\mu_\lambda,f)\bigr)
  =\frac{m'}{m}(\mu,\lambda)\operatorname{tr}I(\mu_{S,\lambda},f_S)+
  \operatorname{tr}\bigl(R_S^{-1}R_S'(\mu_S,\lambda)I(\mu_{S,\lambda},f_S)
  \bigr).
  \]
\item $m(\mu^w,-\lambda)m(\mu,\lambda)=1$; for $\lambda\in i\R$,
  $m(\mu,\lambda)$ is holomorphic of absolute value $1$, $\frac{m'}m(\mu,
  \lambda)$ is real, and $\frac{m'}m(\mu^w,-\lambda)=\frac{m'}m(\mu,\lambda)$.
\item With $\Lambda^{\mathrm{bad}}(s,\chi^2)=\prod_{v\notin
  S_{\mathrm{bad}}(\mu)}L_v(s,\chi_v^2)$, the product over all places
  outside $S_{\mathrm{bad}}(\mu)$, archimedean ones included,
  \[
  \frac{m'}m(\mu,\lambda)=2\frac{(\Lambda^{\mathrm{bad}})'}
  {\Lambda^{\mathrm{bad}}}(2\lambda,\chi^2)-2\frac{(\Lambda^{\mathrm{bad}})'}
  {\Lambda^{\mathrm{bad}}}(2\lambda+1,\chi^2)+\sum_{v\in S_{\mathrm{bad}}
  (\mu)}\frac{r_v'}{r_v}(\mu_v,\lambda).
  \]
\end{enumerate}
\end{corollary}

\begin{proof}
(a) Lemma~\ref{lem:intertwining}(a)--(b), with $S_1=S_k$, and
$M_v=r_vR_v$; then the Leibniz rule, as for Corollary~\ref{cor:mprime}.
(b) $M(w,-\lambda)M(w,\lambda)=1$ (Theorem~\ref{thm:eisenstein}(c))
and~\ref{N2} give the first statement, and the unitarity of $M(w,\lambda)$
and of $R_S$ on $i\R$ the second; differentiate. (c) is the logarithmic
derivative of~\eqref{eq:global-m}.
\end{proof}

\begin{remarkrev}\label{rem:bad-places}
\emph{The factors at the bad places.} For $v\in S_{\mathrm{bad}}(\mu)$
(dyadic places, and places where $\psi_v$ or $\mu_v$ is ramified) we use
the factors $r_v$ of~\cite{LiII} without making them explicit. The
expected choice is again $L_v(2\lambda,\chi_v^2)/L_v(2\lambda+1,\chi_v^2)$
times an exponential factor, which would give
$m(\mu,\lambda)=c\,q^{-\lambda}\Lambda(2\lambda,\chi^2)/\Lambda(2\lambda+1,
\chi^2)$ with the completed Hecke $L$-function and constants $c$, $q>0$,
as the Shimura correspondence suggests; we have not verified~\ref{N2}
for this choice at those places. The distributions $I_A(\gamma,f)$ and the function $m'/m$
below depend on the choice of the $r_v$; the invariant trace formula
holds for each choice. For local coefficients and Plancherel measures of
the genuine principal series of covers of $SL(2,k_v)$ see
Szpruch~\cite{Sz09} and Goldberg--Szpruch~\cite{GS16}; we have not derived
the factors $r_v$ from those results.
\end{remarkrev}

\subsection{The map \texorpdfstring{$\varphi_A$}{phi\_A} and the
distributions \texorpdfstring{$I_A(\gamma,f)$}{I\_A(gamma,f)}}
\label{sec:IAgamma}

We now define the invariant distributions. Fix normalizing factors as in
Section~\ref{sec:normalized}, and for $\pi\in\Pi_u(\AS)$ let
$R_S(\pi,\lambda)=\bigotimes_{v\in S_k}R_v(\pi_v,\lambda)$. For
$f\in C_c^\infty(\GS)_-$ put
\begin{equation}\label{eq:phiA}
\varphi_A(f,\pi)=-\operatorname{tr}\Bigl(R_S(\pi,0)^{-1}\,\frac{dR_S}
{d\lambda}(\pi,0)\,I(\pi,f)\Bigr),\qquad\pi\in\Pi_u(\AS).
\end{equation}
By Theorem~\ref{thm:normalization}, $\varphi_A(f,\pi_\lambda)=
-\operatorname{tr}(R_S^{-1}R_S'(\pi,\lambda)I(\pi_\lambda,f))$ for
$\lambda\in i\R$, and $R_S^{-1}R_S'=\sum_{v\in S_k}R_v^{-1}R_v'$. This is
Arthur's map $\phi_M(f,\pi)=\operatorname{tr}(\mathcal R_M(\pi,P)I_P(\pi,
f))$~\cite[\S7]{A6}, \cite{A9} for $M=A$, $P=B$: the $(G,A)$-family is
$\mathcal R_B(\Lambda)=1$, $\mathcal R_{\bar B}(\Lambda)=R_S(\pi,0)^{-1}
R_S(\pi,\Lambda)$, and, as for~\eqref{eq:MT},
\[
\mathcal R_A(\pi,B)=\lim_{\Lambda\to0}\Bigl(\frac{1}{\theta_B(\Lambda)}+
\frac{R_S(\pi,0)^{-1}R_S(\pi,\Lambda)}{\theta_{\bar B}(\Lambda)}\Bigr)
=-R_S(\pi,0)^{-1}R_S'(\pi,0),
\]
with $\theta_B(\Lambda)=\Lambda=-\theta_{\bar B}(\Lambda)$ in the
normalization of Remark~\ref{rem:alpha}.

\begin{proposition}\label{prop:phiA}
Let $f\in C_c^\infty(\GS)_-$ and $y\in\GS$.
\begin{enumerate}[label=(\alph*)]
\item $\varphi_A(f,\cdot)$ is continuous and integrable on $\Pi_u(\AS)$.
  More precisely, it vanishes unless the conductors of the $\pi_v$, $v$
  finite, are bounded in terms of $f$, and it decreases rapidly in the
  continuous parameters of $\pi$.
\item $\varphi_A(f,\pi^w)=\varphi_A(f,\pi)$.
\item $\varphi_A(f^y,\pi)=\varphi_A(f,\pi)+\widehat{f_{B,y}}(\pi)+
  \widehat{f_{B,y}}(\pi^w)$.
\item If $v\notin S_k$ is as in~\ref{N4} and $S'=S\cup\{v\}$, then
  $\varphi_A(f\otimes e_v,\pi\otimes\pi_v)=\varphi_A(f,\pi)$ if $\pi_v$ is
  trivial on $A(\OO_v)^*$, and $=0$ otherwise.
\end{enumerate}
\end{proposition}

\begin{proof}
(a) $f$ is bi-invariant under an open subgroup at the finite places, which
bounds the conductors. $\pi\mapsto I(\pi,f)$ is the operator-valued
Fourier transform of a compactly supported smooth function, so its trace
norm on each $\KS$-type decreases rapidly in $\pi$ and in the
$\KS$-type; by~\ref{N2} and~\ref{N3}, $R_S^{-1}R_S'$ is continuous in
$\pi$ and of polynomial growth. (b) By~\ref{N2}, $R_S(\pi^w,\nu)=R_S(\pi,-\nu)^{-1}$, so
$R_S(\pi^w,\nu)^{-1}\frac{d}{d\nu}R_S(\pi^w,\nu)$ at $\nu=0$ is
$R_S'R_S^{-1}$ at $(\pi,0)$, and
$\operatorname{tr}(R_S'R_S^{-1}I(\pi^w,f))=\operatorname{tr}(R_S^{-1}R_S'
R_S^{-1}I(\pi^w,f)R_S)=\operatorname{tr}(R_S^{-1}R_S'I(\pi,f))$
by~\ref{N1}. (c) is~\eqref{eq:trace-variance} with $\mathcal
M=R_S(\pi,\lambda)$ at $\lambda=0$. (d) $I(\pi_v,e_v)$ is the projection
onto $\C\varphi^0_v$ (or $0$), and $R_v^{-1}R_v'\varphi_v^0=0$
by~\ref{N4}.
\end{proof}

For $\tilde a\in\AS$ put
\begin{equation}\label{eq:PhiA}
\Phi_A(\tilde a,f)=\int_{\Pi_u(\AS)}\varphi_A(f,\pi)\,\pi(\tilde a)^{-1}\,
d\pi,
\end{equation}
the function on $\AS$ whose Fourier transform is $\varphi_A(f,\cdot)$,
and for $\gamma\in\AS$ regular and $z\in\{\pm1\}$ define
\begin{equation}\label{eq:IAgamma}
I_A(\gamma,f)=J_A(\gamma,f)-\Phi_A(\gamma,f),\qquad
I_A(z,f)=J_A(z,f)-\Phi_A(z,f).
\end{equation}

\begin{theorem}\label{thm:IA}
Let $f\in C_c^\infty(\GS)_-$.
\begin{enumerate}[label=(\alph*)]
\item $\Phi_A(\cdot,f)$ is a continuous, bounded, anti-genuine function
  on $\AS$, smooth in the archimedean coordinates and locally constant in
  the others, with
  $\Phi_A(w\tilde aw^{-1},f)=\Phi_A(\tilde a,f)$. If $f=\bigotimes_vf_v$,
  its restriction to each set $\{\tilde a\in\AS\mid H_S(\tilde a)=X\}$ is
  compactly supported. It is not compactly supported on $\AS$ in general.
\item The distributions $I_A(\gamma,\cdot)$ and $I_A(z,\cdot)$ are
  invariant: $I_A(\gamma,f^y)=I_A(\gamma,f)$ for all $y\in\GS$.
\item $I_A(w\gamma w^{-1},f)=I_A(\gamma,f)$.
\item If $v\notin S_k$ is as in~\ref{N4}, $S'=S\cup\{v\}$, and
  $\gamma_v\in A(\OO_v)^*$ lies over $a(t)$ with $|t^2-1|_v=1$, then
  $I_A(\gamma\gamma_v,f\otimes e_v)=I_A(\gamma,f)$.
\end{enumerate}
\end{theorem}

\begin{proof}
(a) Continuity and boundedness follow from
Proposition~\ref{prop:phiA}(a), smoothness from the rapid decrease, local
constancy from the bound on the conductors, and
the symmetry from Proposition~\ref{prop:phiA}(b) and the invariance of
$d\pi$ under $\pi\mapsto\pi^w$. For $f=\bigotimes f_v$, by the Leibniz
rule $\varphi_A(f,\pi)=\sum_{v\in S_k}\varphi_{A,v}(f_v,\pi_v)\prod_{u\ne
v}\operatorname{tr}I(\pi_u,f_u)$, and $\operatorname{tr}I(\pi_u,f_u)=
\widehat{F_u}(\pi_u)$ with $F_u(\tilde a)=\delta_B(\tilde a)^{1/2}\int
\int f_u(\tilde k^{-1}\tilde an\tilde k)\,dn\,d\tilde k$ compactly
supported on $\overline{A(k_u)}$. So $\Phi_A=\sum_v\Phi_{A,v}\otimes
\bigotimes_{u\ne v}F_u$ with $\Phi_{A,v}$ continuous. On
$\{H_S=X\}$ the $v$-th summand is supported where $\tilde a_u\in
\mathrm{supp}\,F_u$ for $u\ne v$, and then $H_v(\tilde a_v)=X-\sum_{u\ne
v}H_u(\tilde a_u)$ is bounded; this set is compact. The last statement is
Example~\ref{ex:Phi-real}.

(b) By Proposition~\ref{prop:phiA}(c) and Fourier inversion for
$f_{B,y}\in C_c^\infty(\AS)_-$ (Lemma~\ref{lem:uprime}(b)),
\[
\Phi_A(\tilde a,f^y)=\Phi_A(\tilde a,f)+f_{B,y}(\tilde a)+f_{B,y}(w\tilde
aw^{-1}),
\]
since the inverse transform of $\pi\mapsto\hat h(\pi^w)$ is $\tilde
a\mapsto h(w\tilde aw^{-1})$. Subtract this from
Lemma~\ref{lem:variance-local}(b), and at $\tilde a=z$ from
Lemma~\ref{lem:variance-local}(c).

(c) $J_A(w\gamma w^{-1},f)=J_A(\gamma,f)$ by the substitution $x\mapsto
wx$ in~\eqref{eq:JA}, as after Theorem~\ref{thm:fine}; and~(a). (d) For
such $\gamma_v$ the orbital integral of $e_v$ is $1$ and its weighted
orbital integral is $0$, because $v_{A,v}=0$ on the support (the proof of
Theorem~\ref{thm:fine}); and Proposition~\ref{prop:phiA}(d) gives
$\Phi_A(\gamma\gamma_v,f\otimes e_v)=\Phi_A(\gamma,f)$.
\end{proof}

For connected reductive groups, $I_M(\gamma,f)$ is defined in~\cite[\S2]{A9}
by $I_M(\gamma,f)=J_M(\gamma,f)-\sum_{L\in\mathcal L(M),\,L\ne G}
\hat I^L_M(\gamma,\phi_L(f))$, and for covers in~\cite[\S4.5]{LiIV}. Here
$\mathcal L(A)=\{A,G\}$, the only term is $L=A$, and
$\hat I_A^A(\gamma,\phi_A(f))$ is the value at $\gamma$ of the function
with Fourier transform $\phi_A(f)$, which is $\Phi_A(\gamma,f)$. For a
general group the corresponding step needs the trace Paley--Wiener
theorem for $L$; for $L=A$ it is Fourier inversion on an abelian group.
Two further remarks.

\emph{The space containing $\varphi_A(f)$.} Arthur's
$\phi_M$ maps $\mathcal H(G(F_S))$ to a space $\mathcal I_{\mathrm{ac}}
(M(F_S))$ of functions $\phi$ on $\Pi_{\mathrm{temp}}(M(F_S))$ for which
\[
\phi(\pi,X)=\int_{i\mathfrak a_M^*}\phi(\pi_\Lambda)\,e^{-\Lambda(X)}\,
d\Lambda,\qquad X\in\mathfrak a_M,
\]
is, for each $X$, a Paley--Wiener function of $\pi$~\cite[\S11]{A11}.
In our
case $\phi(\pi,X)$ is the Fourier transform of the restriction of
$\Phi_A(\cdot,f)$ to $\{H_S=X\}$, and Theorem~\ref{thm:IA}(a) is the
statement that $\varphi_A(f)\in\mathcal I_{\mathrm{ac}}$. This space is
strictly larger than the space $\mathcal I(\overline{A(\A_S)})$ of Fourier
transforms of compactly supported functions, for the reason given in
Section~\ref{sec:local-dist} for the spaces of test functions:
$\mathfrak a_A\ne0$ (Example~\ref{ex:Phi-real}).

\emph{Weighted characters.} The weighted character of~\cite{A9} is
$J_A(\pi,f)=\operatorname{tr}(\mathcal R_A(\pi,B)I(\pi,f))=\varphi_A(f,
\pi)$, so the invariant distribution $I_A(\pi,X,f)=J_A(\pi,X,f)-
\hat I^A_A(\pi,X,\varphi_A(f))$ vanishes for every $\pi\in\Pi_u(\AS)$.
This does not
make the spectral side of the invariant formula vanish: the terms that
remain are characters, with the coefficients of Theorem~\ref{thm:itf}.

\begin{example}\label{ex:Phi-real}
Let $v$ be real, $\pi_v$ and $m_0$ as in Lemma~\ref{lem:explicit-r}(b),
and let $P_j$ be the projection of $V(\pi_v)$ onto the weight $m_0+2j$.
For $f_v\in C_c^\infty(\overline{SL(2,\R)})_-$ let $F_j$ be the function
on $\overline{A(\R)}$ with $\widehat{F_j}(\pi)=\operatorname{tr}(P_jI(\pi,
f_v))$; then $\sum_jF_j=F_{f_v}$, the Abel transform in the proof of
Theorem~\ref{thm:IA}(a). By~\eqref{eq:rho-m}, $\varphi_{A,v}(f_v,\pi_{v,it})
=\sum_j\rho_j(t)\widehat{F_j}(\pi_{v,it})$, and since
$\frac1{2\pi}\int_\R\frac{2c}{c^2+t^2}e^{-itX}\,dt=e^{-c|X|}$ for $c>0$,
\[
\Phi_{A,v}(\tilde a,f_v)=\sum_j\int_\R\kappa_j(X)\,F_j\bigl(\tilde a\,a(e^X)
\bigr)\,dX,\qquad\kappa_j(X)=\sum_ie^{-c_i|X|},
\]
with $c_i$ as in~\eqref{eq:rho-m}. So $\Phi_{A,v}(\cdot,f_v)$ is the Abel
transform smoothed, on each $\overline{K_v}$-type, by a kernel
$\kappa_j\ge0$ that decays exponentially but is not compactly supported;
it vanishes on the weight $m_0$. For example, if $f_v$ transforms on
both sides under $\overline{K_v}$ by the character of weight $m_0+2j$
with $j\ne0$, and $F_j=F_{f_v}\ge0$ is not zero, then $\Phi_{A,v}(\tilde
a,f_v)>0$ at every $\tilde a$ in the cosets of $\{a(s)\mid s>0\}$ that
meet the support of $F_{f_v}$, although $F_{f_v}$ is compactly supported.
\end{example}

\subsection{The invariant trace formula}\label{sec:itf}

We return to $f=f_S\otimes e^S$ with $f_S\in\mathcal H(\GS)$, and to the
situation of Theorem~\ref{thm:tf}. For $\gamma\in A(\Q)$ put
$\Phi_A(\gamma,f)=\Phi_A(\gamma_S,f_S)$ if $\gamma\in K_v$ for all
$v\notin S_k$, and $\Phi_A(\gamma,f)=0$ otherwise; for
$\mu\in\Pi_{\mathrm{aut}}(X_A)$ put $\varphi_A(f,\mu_\lambda)=\varphi_A(f_S,
\mu_{S,\lambda})$ if $\mu_v$ is trivial on $A(\OO_v)^*$ for all
$v\notin S_k$, and $0$ otherwise.

\begin{lemma}\label{lem:poisson-phi}
$\Phi_A(\gamma,f)=0$ for all but finitely many $\gamma\in A(\Q)$, and
\[
\sum_{\mu\in\Pi_{\mathrm{aut}}(X_A)}\frac1{2\pi}\int_\R\varphi_A(f,\mu_{it})
\,dt=c_k\sum_{\gamma\in A(\Q)}\Phi_A(\gamma,f),
\]
the left side converging absolutely.
\end{lemma}

\begin{proof}
We may take $f_S=\bigotimes_vf_v$. Let $h=\Phi_A(\cdot,f_S)\otimes
e_A^S$ on $\overline{A(\A)}$, where $e_A^S=\bigotimes_{v\notin S_k}e_{A,v}$
is the product of the units of Lemma~\ref{lem:transfer}(b). By
Theorem~\ref{thm:IA}(a) the restriction of $h$ to $\overline{A(\A)^1}$ is
compactly supported; it is smooth at the archimedean places and locally
constant at the finite ones, because $\varphi_A(f_S,\cdot)$ decreases
rapidly and is supported on finitely many components. So
$\gamma\mapsto h(\gamma)=\Phi_A(\gamma,f)$ has finite support on the
discrete subgroup $A(\Q)$, and Theorem~\ref{thm:tf-rank0} applies to
$h|_{\overline{A(\A)^1}}$: $\sum_\mu\int_{\overline{A(\A)^1}}h\mu=c_k
\sum_\gamma h(\gamma)$. Finally $\int_{\overline{A(\A)^1}}h\mu=\frac1{2\pi}
\int_\R\hat h(\mu_{it})\,dt$ and $\hat h(\mu_{it})=\varphi_A(f,\mu_{it})$,
by Fourier inversion on $\R$. The left side converges absolutely by
Proposition~\ref{prop:phiA}(a): the $\mu$ that contribute form finitely
many cosets of a lattice in the unitary dual of $A(\Q_S)$ modulo
$i\R$.
\end{proof}

For $\mathcal O\in\mathcal C$ and $\chi\in\mathcal X$ define
\begin{equation}\label{eq:Iochi}
I_{\mathcal O}(f)=J_{\mathcal O}(f)-\frac{c_k}2\sum_{\gamma\in\mathcal O\cap
A(\Q)}\Phi_A(\gamma,f),\qquad
I_\chi(f)=J_\chi(f)-\frac12\sum_{\mu\in\chi}\frac1{2\pi}\int_\R
\varphi_A(f,\mu_{it})\,dt,
\end{equation}
the second sum being empty for $\chi\in\mathcal X_G$. That is, $I=J-J'$
with
$J'=|W_G|^{-1}(I^A)^\wedge(\varphi_A(f))$, where $(I^A_{\mathcal O})^\wedge$
and $(I^A_\chi)^\wedge$ are the distributions~\eqref{eq:JAo} transported
to the Fourier transform side.

\begin{theorem}[The invariant trace formula]\label{thm:itf}
Let $S$ and $f=f_S\otimes e^S$ be as in Theorem~\ref{thm:tf}. Then
\[
\sum_{\mathcal O\in\mathcal C}I_{\mathcal O}(f)=\sum_{\chi\in\mathcal X}
I_\chi(f),
\]
both sides converging absolutely, and:
\begin{enumerate}[label=(\alph*)]
\item for the central classes $\mathcal O_z$, $z\in\{\pm1\}$,
  \[
  I_{\mathcal O_z}(f)=a^G(S,1)f_S(z)+\sum_{\xi\in k_S^\times/k_S^{\times2}}
  a^G(S,u_\xi)\,J_G(zu_\xi,f_S)+\tfrac12\,c_k\,I_A(z,f_S);
  \]
\item for $\mathcal O$ elliptic, $I_{\mathcal O}(f)=J_{\mathcal O}(f)=
  \mathrm{vol}(T_\mu(\Q)\backslash T_\mu(\A))\,J_G(\mu_S,f_S)$;
\item for $\mathcal O$ hyperbolic, the class of $\gamma=a(t)$,
  $I_{\mathcal O}(f)=c_k\,I_A(\gamma_S,f_S)$ if $t\in\OO_v^\times$ for all
  $v\notin S_k$, and $I_{\mathcal O}(f)=0$ otherwise;
\item for $\chi$ the class of $(G,\pi)$, $I_\chi(f)=m(\pi)\operatorname{tr}
  \pi(f)$;
\item for $\chi\in\mathcal X_A$ unramified and $\mu\in\chi$,
  \[
  I_\chi(f)=-\frac1{2\pi}\int_\R\frac{m'}m(\mu,it)\,\operatorname{tr}
  I(\mu_{it},f)\,dt;
  \]
\item for $\chi=\chi_b$ ramified and $\mu=\chi_{\psi_b}$,
  \[
  I_\chi(f)=\operatorname{tr}\bigl(R(f)|_{V_b}\bigr)-\frac12\cdot\frac1{2\pi}
  \int_\R\frac{m'}m(\mu,it)\,\operatorname{tr}I(\mu_{it},f)\,dt
  +\frac14\operatorname{tr}\bigl(M(w,0)\,I(\mu_0,f)\bigr).
  \]
\end{enumerate}
Each of the distributions (a)--(f) is the restriction to $\mathcal H(\GS)$
of an invariant distribution in $f_S\in C_c^\infty(\GS)_-$. Here
$a^G(S,u_\xi)$ are the coefficients of Theorem~\ref{thm:fine}, and $m$ and
$I_A$ are formed with the normalizing factors fixed in
Section~\ref{sec:normalized}.
\end{theorem}

\begin{proof}
By Theorem~\ref{thm:tf} and Lemma~\ref{lem:poisson-phi}, $\sum_{\mathcal
O}I_{\mathcal O}(f)-\sum_\chi I_\chi(f)=-\frac12\bigl(c_k\sum_\gamma
\Phi_A(\gamma,f)-\sum_\mu\frac1{2\pi}\int\varphi_A(f,\mu_{it})\,dt\bigr)=0$.
(a) is Theorem~\ref{thm:tf}(a) and~\eqref{eq:IAgamma}, since $\mathcal O_z
\cap A(\Q)=\{z\}$. (b) $\mathcal O\cap A(\Q)$ is empty. (c) $\mathcal
O\cap A(\Q)=\{\gamma,\gamma^{-1}\}$ and $\Phi_A(\gamma^{-1},f)=\Phi_A(\gamma,
f)$ by Theorem~\ref{thm:IA}(a), because $\gamma^{-1}=w\gamma w^{-1}$; with
Theorem~\ref{thm:tf}(c) this gives $c_k(J_A-\Phi_A)(\gamma_S,f_S)$. If $t$
is not a unit at some $v\notin S_k$, both $J_{\mathcal O}(f)$ and
$\Phi_A(\gamma,f)$ vanish. (d) is Theorem~\ref{thm:tf}(d). (e), (f) By
Theorem~\ref{thm:Jchi} at $T=0$ and Corollary~\ref{cor:global-m}(a),
\[
-\operatorname{tr}\bigl(M^{-1}M'(w,it)I(\mu_{it},f)\bigr)=-\frac{m'}m(\mu,
it)\operatorname{tr}I(\mu_{it},f)+\varphi_A(f,\mu_{it})
\]
when $\mu$ is unramified outside $S$, and both sides vanish otherwise. In
the unramified case the two members of $\chi$ contribute equally
to~\eqref{eq:Iochi}, by Proposition~\ref{prop:phiA}(b) and $t\mapsto-t$,
so $\varphi_A$ cancels; in the ramified case $J_\chi$ has the factor
$\frac12$ and $\chi$ has one member. Invariance: (b) and (d) are an
orbital integral and a character; the terms of~(a) and~(c) are invariant
by Theorem~\ref{thm:IA}(b), and $f_S(z)$ and $J_G(zu_\xi,\cdot)$ are
orbital integrals. In (e) and (f), $\operatorname{tr}I(\mu_{it},\cdot)$
and $\operatorname{tr}(R(\cdot)|_{V_b})$ are characters, and
$\operatorname{tr}(M(w,0)I(\mu_0,\cdot))$ is invariant because $M(w,0)$
commutes with $I(\mu_0)$. The absolute convergence is that of
Theorem~\ref{thm:tf} and Lemma~\ref{lem:poisson-phi}.
\end{proof}

Theorem~\ref{thm:itf} is the invariant trace formula of~\cite{LiIV} for
this group. Its terms are explicit with two exceptions, which are
defined but not evaluated here. The unipotent coefficients
$a^G(S,u_\xi)$ in~(a) are those of Theorem~\ref{thm:fine}
(Remark~\ref{rem:sec4-status}). And the normalizing factors $r_v$ at the
places of $S_{\mathrm{bad}}(\mu)$, which always include the dyadic places,
are those of Theorem~\ref{thm:normalization}; the function $m'/m$
in~(e) and~(f) and the distributions $I_A$ in~(a) and~(c) depend on them
(Remark~\ref{rem:bad-places}). The elliptic terms~(b), the discrete
terms~(d), the residual term and the ramified term in~(f) do not involve
either.

It remains to evaluate $M(w,0)$ in~(f). For $a\in k^\times$ let $\eta_a$
be the quadratic character $t\mapsto\prod_v(-a,t_v)_v$ of
$k^\times\backslash\A_k^\times$.

\begin{lemma}\label{lem:whittaker}
Let $\mu=\chi_\psi\cdot(\chi\circ p)\in\Pi_{\mathrm{aut}}(X_A)$,
$\varphi=\bigotimes_v\varphi_v\in\mathcal H_B(\mu)$ and $a\in k^\times$.
For $\mathrm{Re}\,\lambda>1$ put
\[
W_a(x,\varphi,\lambda)=\int_{k\backslash\A_k}E\bigl(n(u)x,\varphi,\lambda
\bigr)\,\psi(-au)\,du .
\]
\begin{enumerate}[label=(\alph*)]
\item $E(x,\varphi,\lambda)=E_B(x,\varphi,\lambda)+\sum_{a\in k^\times}
  W_a(x,\varphi,\lambda)$, and
  $W_a(x,\varphi,\lambda)=\prod_vW_{a,v}(x_v,\varphi_v,\lambda)$ with
  $W_{a,v}(x_v,\varphi_v,\lambda)=\int_{k_v}\varphi_{v,\lambda}\bigl((w,1)^{-1}
  (n(u),1)x_v\bigr)\psi_v(-au)\,du$.
\item Each $W_{a,v}(x_v,\varphi_v,\lambda)$ is an entire function of
  $\lambda$.
\item If $v$ is as in~\ref{N4}, $|a|_v=1$, $\varphi_v=\varphi_v^0$ and
  $x_v\in K_v^*$, then
  \[
  W_{a,v}(x_v,\varphi^0_v,\lambda)=1+\eta_{a,v}(\varpi)\chi_v(\varpi)
  q^{-\lambda-\frac12}=\frac{L_v(\lambda+\frac12,\chi_v\eta_{a,v})}
  {L_v(2\lambda+1,\chi_v^2)} .
  \]
\item Hence, for a finite set $S_1$ of places containing those not covered
  by~(c),
  \[
  W_a(x,\varphi,\lambda)=\frac{L^{S_1}(\lambda+\frac12,\chi\eta_a)}
  {L^{S_1}(2\lambda+1,\chi^2)}\prod_{v\in S_1}W_{a,v}(x_v,\varphi_v,\lambda),
  \]
  which continues $W_a$ meromorphically to $\C$.
\end{enumerate}
\end{lemma}

\begin{proof}
(a) The first formula is the Fourier expansion on $k\backslash\A_k$. By
the Bruhat decomposition $G(\Q)=B(\Q)\sqcup B(\Q)w^{-1}N(\Q)$, for
$a\ne0$ the series~\eqref{eq:Eis} unfolds to $\int_{\A_k}\varphi_\lambda
(w^{-1}n(u)x)\psi(-au)\,du$, which is the product. (b) For $v$ finite and
$|u|$ large, $w^{-1}n(u)=n(-u^{-1})a(-u^{-1})\bar n(u^{-1})$ with $\bar
n(u^{-1})$ in a subgroup that fixes $I(x_v)\varphi_v$, so the integrand
is $|u|^{-\lambda-1}$ times a function of $u$ that is invariant under
multiplication by a subgroup $1+\varpi^c\OO$; its integral against
$\psi_v(-au)$ over a shell $|u|=q^m$ vanishes for $m$ large, and
$W_{a,v}$ is a polynomial in $q^{\pm\lambda}$. For $v$ archimedean the
integrand is smooth in $u$ with an asymptotic expansion in powers of
$u^{-1}$ times $|u|^{-\lambda-1}$, and repeated integration by parts
against $\psi_v(-au)$ continues the integral to all $\lambda$. (c)
Proceed as for Lemma~\ref{lem:intertwining}(b). On $|u|\le1$ the
integrand is $1$. On the shell $|u|=q^m$, $m\ge1$, write
$u=\varpi^{-m}u_0$; the decomposition $(w,1)^{-1}(n(u),1)=(n(-u^{-1}),1)
(a(-u^{-1}),1)(\bar n(u^{-1}),1)$ holds in the cover, because every
value of Kubota's cocycle that occurs is of the form $(1,\cdot)_v$ or
$(u,-u)_v$, hence $1$, and $\kappa_v=1$ on all factors. So the integrand is
$\mu_v(a(-u^{-1}),1)|u|^{-\lambda-1}=\chi_v(\varpi)^m|u|^{-\lambda-1}
\gamma_v(-\varpi^mu_0)$, and by Lemma~\ref{lem:weil}(a)--(b)
$\gamma_v(-\varpi^mu_0)$ is $1$ for $m$ even and $(-u_0,\varpi)_v
\gamma_v(\varpi)$ for $m$ odd. For $m$ even the integral of $\psi_v(-au)$
over the shell is $0$ if $m\ge2$. For $m$ odd the integral of
$(u_0,\varpi)_v\psi_v(-au)$ is $0$ if $m\ge3$, and for $m=1$ it is the
Gauss sum $\sum_{u_0\in(\OO/\varpi)^\times}(u_0,\varpi)_v\psi_v(-a\varpi^{-1}
u_0)=(-a,\varpi)_v\,q^{1/2}\gamma_v(\varpi)$, by~\eqref{eq:weil-index}. The
shell $m=1$ therefore contributes $(-1,\varpi)_v(-a,\varpi)_v\gamma_v
(\varpi)^2\chi_v(\varpi)q^{-\lambda-\frac12}$, and
$\gamma_v(\varpi)^2=(-1,\varpi)_v$. (d) follows.
\end{proof}

Part~(c) is the unramified Whittaker function of a genuine principal
series at the identity; compare~\cite{GPS80}. For $SL(2)$ the same
computation gives $1-\chi_v(\varpi)q^{-\lambda-1}=L_v(\lambda+1,\chi_v)
^{-1}$. The script checks both at $p=3,5,7,11$.

\begin{proposition}\label{prop:Mw0}
Let $b\in k^\times$ and $\mu=\chi_{\psi_b}$. Then $E(x,\varphi,0)=0$ for
all $\varphi\in\mathcal H_B(\mu)$, and
\[
M(w,0)=-1\quad\text{on }\mathcal H_B(\chi_{\psi_b}).
\]
Consequently the ramified term of Theorem~\ref{thm:Jchi}(c) is
$\frac14\operatorname{tr}(M(w,0)I(\mu_0,f))=-\frac14\operatorname{tr}
I(\chi_{\psi_b,0},f)$.
\end{proposition}

\begin{proof}
Write $\mu=\chi_\psi\cdot(\chi\circ p)$ with $\chi=\prod_v(b,\cdot)_v$
(Lemma~\ref{lem:weil}(d)), so $\chi^2=1$. It suffices to treat
factorizable $\varphi$. By
Theorem~\ref{thm:eisenstein}(d), $E(x,\varphi,\lambda)$ and
$M(w,\lambda)$ are holomorphic at $\lambda=0$.

\emph{Step 1: the Whittaker coefficients vanish at $\lambda=0$.} In
Lemma~\ref{lem:whittaker}(d) the denominator is
$L^{S_1}(2\lambda+1,\chi^2)=\zeta_k^{S_1}(2\lambda+1)$, because $\chi^2=1$.
The Dedekind zeta function has a pole at $s=1$, and removing finitely
many Euler factors does not remove it; so
$\zeta_k^{S_1}(2\lambda+1)^{-1}$ has a zero at $\lambda=0$. The numerator
$L^{S_1}(\lambda+\frac12,\chi\eta_a)$ is the $L$-function of a quadratic
Hecke character at the center of the critical strip, which is finite
(for $\chi\eta_a=1$ it is $\zeta_k^{S_1}(\frac12)$), and the local factors
at $S_1$ are entire. So $W_a(x,\varphi,0)=0$ for every $a\in k^\times$ and
every $x$. The zero is specific to the cover: the unramified Whittaker
function has the denominator $L_v(2\lambda+1,\chi_v^2)$, in which the
character is squared, so every ramified class sees the pole of
$\zeta_k$.

\emph{Step 2: the Eisenstein series is its constant term.} By
Lemma~\ref{lem:whittaker}(a) and Theorem~\ref{thm:eisenstein}(b),
\[
F(x):=E(x,\varphi,0)=E_B(x,\varphi,0)=\bigl(\varphi+M(w,0)\varphi\bigr)(x)
\,e^{H(x)}.
\]

\emph{Step 3: $|F|$ is constant.} $F$ is left invariant under $G(\Q)$.
Being equal to its constant term $\int_{N(\Q)\backslash N(\A)}F(nx)\,dn$,
it is left invariant under $N(\A)$. Since $w\in G(\Q)$, it is then left
invariant under $wN(\A)w^{-1}$ as well: $F(wnw^{-1}x)=F(nw^{-1}x)=
F(w^{-1}x)=F(x)$. Here $wN(\A)w^{-1}$ is the image of the lower triangular
unipotent group $\bar N(\A)$ under its canonical section
(Lemma~\ref{lem:local}(h)). Now $N(\A)$ and $\bar N(\A)$ generate
$G(\A)$: over a field or a local ring every element of $SL(2)$ is a
product of a bounded number of elementary matrices, the bound being
independent of the ring, so the factorization can be carried out at all
places simultaneously, with integral factors at almost all places. Hence
the subgroup of $\GA$ generated by the lifts of $N(\A)$ and $\bar N(\A)$
maps onto $G(\A)$, and $F$ is left invariant under it. Since $F$ is
genuine, $|F|$ is invariant under $\mu_2$ too. So $|F|$ is left invariant
under all of $\GA$, that is, constant.

\emph{Step 4: the constant is $0$.} For $\tilde a\in\overline{A(\A)}$ the
formula of Step~2 gives $|F(\tilde ax)|=e^{H(\tilde a)}|F(x)|$, because
$\varphi$ and $M(w,0)\varphi$ both lie in $\mathcal H_B(\mu)$ when
$\mu^w=\mu$. Choosing $H(\tilde a)\ne0$ we get $F=0$, and restricting to
$\overline{K_0}$ gives $\varphi+M(w,0)\varphi=0$.
\end{proof}

The statement is about representations, not only about $L$-functions:
$M(w,0)$ is a self-intertwining operator of the unitary representation
$I(\mu_0)$ with $M(w,0)^2=1$, and the proposition says that it is the
scalar $-1$ for every ramified class. The proof finds the sign from a
global fact: an automorphic function that is equal to its constant term
and homogeneous of degree one under $A$ must vanish.

For $SL(2)$ the same argument applies to the trivial character, where
$L^{S_1}(\lambda+1,1)^{-1}$ vanishes at $\lambda=0$: this is the classical
fact that the scattering matrix of $SL(2,\Z)$ is $-1$ at $s=\frac12$ and
$E(z,\frac12)=0$. For a quadratic character $\eta\ne1$ of $SL(2)$ it does
not apply, and $M(w,0)$ is then a product of nonscalar local
operators~\cite{LL}; that is the source of the multiplicity formula
in~\cite{LL}. On the double cover the denominator is always
$\zeta_k^{S_1}(2\lambda+1)$, and every ramified class behaves like the
trivial character of $SL(2)$. This agrees with the irreducibility of the
unitary principal series of $\overline{SL(2,k_v)}$~\cite{Mo2}. For
$k=\Q$ and $b\in\{\pm1,\pm2\}$ the script confirms $M(w,0)=-1$ by a
separate computation, from the residues at $\lambda=0$ of the real and
the $2$-adic local operators and the zero of order $|S|$ of
$m^S(\mu,\lambda)$, $S=\{\infty,2\}$.

\begin{corollary}\label{cor:itf-final}
In Theorem~\ref{thm:itf},
\begin{multline*}
\sum_{\chi\in\mathcal X}I_\chi(f)=\sum_{\pi}m_{\mathrm{disc}}(\pi)
\operatorname{tr}\pi(f)-\frac14\sum_{b\in k^\times/k^{\times2}}
\operatorname{tr}I(\chi_{\psi_b,0},f)\\
-\frac12\sum_{\mu\in\Pi_{\mathrm{aut}}(X_A)}\frac1{2\pi}\int_\R\frac{m'}m
(\mu,it)\operatorname{tr}I(\mu_{it},f)\,dt,
\end{multline*}
where $\pi$ runs over the irreducible constituents of
$L^2_{\mathrm{disc}}=L^2_{\mathrm{cusp}}\oplus\bigoplus_bV_b$ with their
multiplicities $m_{\mathrm{disc}}(\pi)$. Every term is a character of
$\GS$ evaluated at $f_S$.
\end{corollary}

\begin{proof}
Theorem~\ref{thm:itf}(d)--(f), Proposition~\ref{prop:Mw0}, and
Corollary~\ref{cor:global-m}(b) to write~(e) as half the sum over the two
members of $\chi$.
\end{proof}

In Arthur's notation~\cite[\S4]{A10} this is $\sum_t\int_{\Pi(G,t)}a^G(\pi)
I_G(\pi,f)\,d\pi$: the coefficient $a^G(\pi)$ is $m_{\mathrm{disc}}(\pi)$ on
the discrete spectrum, $-\frac14$ at $I(\chi_{\psi_b,0})$, and
$-\frac1{4\pi}\frac{m'}m(\mu,it)\,dt$ on the unitary principal series; the
first two make up the discrete part $I_{\mathrm{disc}}$. Since
$V_b\cong r^+_{\psi_{-b}}$ (Theorem~\ref{thm:residual}(c)),
$\operatorname{tr}(R(f)|_{V_b})=\prod_v\operatorname{tr}r^+_{\psi_{-b},v}
(f_v)$, the character of the even Weil representation.

\begin{example}\label{ex:itf-Q}
Let $k=\Q$ and $S=\{\infty,2\}$, so $f=f_\infty\otimes f_2\otimes e^S$.
By Example~\ref{ex:ramified-Q} the ramified classes that contribute are
$\chi_{\psi_b}$, $b\in\{\pm1,\pm2\}$, and the discrete correction is
$-\frac14\sum_{b=\pm1,\pm2}\operatorname{tr}I(\chi_{\psi_b,0},f)$. For
$\mu=\chi_{\psi_b}$ we have $S_{\mathrm{bad}}(\mu)=\{2\}$ and, with
$\xi(s)=\pi^{-s/2}\Gamma(\frac s2)\zeta(s)$,
\[
m(\mu,\lambda)=c\,\frac{\xi(2\lambda)}{\xi(2\lambda+1)}\cdot
\frac{1-2^{-2\lambda}}{1-2^{-1-2\lambda}}\cdot r_2(\mu_2,\lambda),\qquad
\frac{\xi(2\lambda)}{\xi(2\lambda+1)}\Big|_{\lambda=0}=-1,
\]
where $r_2$ is the $2$-adic factor of Theorem~\ref{thm:normalization}
(Remark~\ref{rem:bad-places}) and $c=\sqrt2\,e^{\pm i\pi/4}$ comes from
Lemma~\ref{lem:explicit-r}(b), the sign being that of $b$. The
hyperbolic classes that contribute are those of $a(t)$, $t=\pm2^j$,
$j\ge1$, and their terms are $c_\Q\,I_A(a(t)_S,f_S)$. The archimedean part of the correction $\Phi_A$ is
given by Example~\ref{ex:Phi-real}.
\end{example}

\subsection{Status of the results}\label{sec:status6}

\begin{remarkrev}\label{rem:sec6-status}
\emph{What rests on citation.} Proved here, given
Sections~\ref{sec:rank0}--\ref{sec:spectral} and
Theorem~\ref{thm:normalization}: Lemmas~\ref{lem:uprime}--%
\ref{lem:conj-identity}; Lemma~\ref{lem:explicit-r}(b), and~(a) except
for rationality; Corollary~\ref{cor:global-m}; Proposition~\ref{prop:phiA};
Theorem~\ref{thm:IA}; Lemma~\ref{lem:poisson-phi}; Theorem~\ref{thm:itf};
Lemma~\ref{lem:whittaker}; Proposition~\ref{prop:Mw0}; and
Corollary~\ref{cor:itf-final}. Taken from the literature:
Theorem~\ref{thm:normalization}~\cite{A11,LiII}, of which only the bad
places are not made explicit (Remark~\ref{rem:bad-places}); the generic
irreducibility of $I(\pi_{v,\lambda})$ and the adjoint formula used in
Lemma~\ref{lem:explicit-r}(a); the estimates behind
Proposition~\ref{prop:phiA}(a) for functions that are not $\KS$-finite;
and, in the proof of Lemma~\ref{lem:variance-global}, the interchange of
integrals on the spectral side. We have not compared the
definition~\eqref{eq:IAgamma} with the definitions of~\cite{LiIV} beyond
the remark after Theorem~\ref{thm:IA}. As in
Theorem~\ref{thm:tf}, the set $S$ must be large for the given $f_S$, and
the coefficients $a^G(S,u_\xi)$ are not computed
(Remark~\ref{rem:sec4-status}).
\end{remarkrev}

\begin{remarkrev}\label{rem:verified-sec6}
The script \texttt{checks/sec6\_n2\_check.py} verifies, for $k=\Q$: the
sign of the cross terms in the Maass--Selberg relation~\eqref{eq:MS},
hence of the ramified term in Theorem~\ref{thm:Jchi}(c), on the
Eisenstein series of $SL(2,\Z)$ (by numerical integration over the
modular surface); Lemma~\ref{lem:uprime}(a) and
Example~\ref{ex:uprime}; Lemma~\ref{lem:variance-local}(a)--(b) at
$v=\R$, with explicit Iwasawa decompositions and by quadrature for a
Gaussian $f$; the identity~\eqref{eq:conj-identity} for the principal
series of $SL(2,\R)$ in the compact picture;
formulas~\eqref{eq:real-M}--\eqref{eq:rho-m} and the Fourier transform in
Example~\ref{ex:Phi-real}; Lemma~\ref{lem:whittaker}(c) at $p=3,5,7,11$,
shell by shell in the cover, and its analog for $SL(2)$;
Proposition~\ref{prop:Mw0} for $b=\pm1,\pm2$, from local residues; and the
value $-1$ in Example~\ref{ex:itf-Q}.
\end{remarkrev}
